\color{unimain}{\subsection{Защита от непереключения фаз}}
\color{black}

\begin{enumerate}

\item
Защита от непереключения фаз (ЗНФ) предназначена для обнаружения расхождения полюсов выключателя с пофазным приводом при его включении или отключении. 

\item
Ввод в работу ЗНФ выключателя В осуществляется программным ключом \textbf{«ЗНФ»}. Факт расхождения полюсов может устанавливаться одним из следующих способов:

\begin{itemize}
	\item на основании внешнего сигнала от сборки блок-контактов выключателя;
	\item функцией фиксации положения выключателя (принцип работы описан в пункте \ref{func:fpv} настоящего руководства по эксплуатации), которая пофазно контролирует состояние блок-контактов выключателя и в случае расхождения полюсов формирует сигнал \textit{«ФПВ В/ ФПВ: Внутр. пуск. ЗНФ»}.
\end{itemize}

\item
Действие на отключение выключателя осуществляется с выдержкой времени \textbf{«Тср ЗНФ»}. Задаваемая выдержка времени предназначена для отстройки от разновременности переключения блок-контактов выключателя.

\item
ЗНФ может быть оперативно выведена из работы от внешнего сигнала \textit{«ЗНФ/ ЗНФ: Вывод»}, назначенного на соответствующий дискретный вход или функциональную клавишу.

\item
Алгоритм работы ЗНФ представлен на рисунке \ref{pict_auv:ZNF}. Параметры настройки функции ЗНФ приведены в таблице \ref{app_set:znf}.

\medskip
\begin{figure}[H]
\centering
\includegraphics[width=0.75\textwidth,height=0.75\textheight,keepaspectratio]{\fbpath/10.02 АУВ/ЗНФ/img/ZNF.pdf}
\captionof{figure}{Функциональная схема ЗНФ}\label{pict_auv:ZNF}
\end{figure}

\end{enumerate}